\documentclass[11pt]{article}
\usepackage[margin=1in]{geometry}
\usepackage{amsmath, amssymb, booktabs, graphicx, hyperref, xcolor}
\usepackage{natbib}
\bibliographystyle{plainnat}
\newcommand{\lamhat}{\hat{\lambda}}
\title{No Eligible Populations: A Pre-Registered Feasibility Census\\ for Trajectory-Level Conformal Risk Control on Public Agent Traces}
\author{Ranadhir Ghosh\\ FIS Global}
\date{}
\begin{document}
\maketitle

\begin{abstract}
Conformal risk control (CRC) can in principle certify an agent
pipeline's trajectory-level failure risk at a pre-set level
$\alpha$, abstaining on trajectories it cannot certify. Whether
that guarantee can be exercised \emph{honestly} on the agent
trajectory corpora the field actually has is an empirical
question that is usually answered by shopping for a corpus that
works. We answer it the other way: we pre-registered a
feasibility census --- frozen data floors (exact Clopper--Pearson
split sizes, a pre-specified base-failure band, task-grouped
disjoint splits), a frozen outcome-binarisation rule that
refuses post-hoc interpretation, and a frozen eligibility
decision table --- over the nine public software-engineering
agent populations that survive a counted-evidence screen
(332{,}499 trajectories in two artefact families). \textbf{No
population is eligible.} All six Claude cells of the SWE-smith
trajectory collection have base failure rates of 0.545--0.610,
above the pre-registered 0.50 ceiling; all three CoderForge
source populations carry an outcome value ($-1.0$, on
0.25--0.84\% of rows) that appears nowhere in the artefact
documentation and that the frozen binarisation may not
interpret. The result is a registered H-DATA outcome: the
binding constraints on trajectory-level CRC today are data-side
and specific --- outcome encodings and base-rate concentration
--- not the guarantee machinery, which we verify controls risk
on planted-signal synthetic traces and which previously ran
end-to-end (to a negative verdict) on the one corpus ever
attempted. We report the full census, the counterfactual rates
that make the cost of the discipline visible, and three
recommendations for trajectory-corpus publishers.
\end{abstract}

\section{Introduction}

An agent that acts over many steps fails or succeeds as a
\emph{trajectory}, and a deployment guarantee worth having is a
statement about trajectories: \emph{among the trajectories this
system accepts, the failure rate is at most $\alpha$} --- with
abstention, not silent risk, as the price of the guarantee.
Conformal risk control \citep{angelopoulos2022conformal}
supplies exactly this shape of statement for a scalar score:
choose a threshold on held-out calibration data so that the
expected loss is controlled at $\alpha$ with a finite-sample,
distribution-free correction. Lifted to agent traces --- a
per-trajectory risk score, a calibration split, an abstention
rule --- it is the natural certification layer for agentic
systems, and learned per-step scorers in the style of process
reward models \citep{lightman2023verify, wang2024mathshepherd,
xi2025agentprm} are the natural scores to certify.

What is missing is not the machinery but the evidence that the
exercise is feasible on real recorded populations under honest
rules. Feasibility here is fragile in a specific way: every
degree of freedom in corpus choice --- which collection, which
subset, what counts as a pass, which base rates are ``in
range'' --- can be, and routinely is, resolved after seeing the
data, at which point the guarantee's honesty is nominal. The
remedy is the one this paper applies: freeze the feasibility
criteria \emph{before} touching outcomes, as a pre-registration,
and let the census fall where it falls.

Our contribution is that census, executed end-to-end under a
frozen registration (itself the third design iteration of a
programme whose earlier pilot and design stages are part of the
record). Over the nine public populations that survived a
prior counting screen on documented evidence, the census finds
\textbf{zero eligible populations}, for two disjoint,
fully documented reasons: base-rate concentration above the
registered ceiling in one family, and an uninterpretable
outcome encoding in the other. Sections 2--4 give the machinery
and the frozen protocol; Section 5 gives the census; Section 6
discusses what the negative result does and does not establish;
Section 7 states recommendations for corpus publishers. The
paper's claim is deliberately narrow: it is a feasibility
result about the present public data landscape for
trajectory-level CRC, produced under rules strict enough that a
positive answer would have meant something.

\section{Background and related work}

\paragraph{Conformal risk control.} CRC
\citep{angelopoulos2022conformal} generalises conformal
prediction from coverage to expected-loss control: for a
monotone loss $L(i, \lambda) = \mathbf{1}\{\text{failed}(i)\}
\cdot \mathbf{1}\{S(i) \le \lambda\}$ in a threshold $\lambda$
on a score $S$, choosing the largest $\lambda$ whose corrected
empirical risk $(\sum_i L(i,\lambda) + 1)/(n+1)$ stays below
$\alpha$ on a calibration sample controls the expected loss at
$\alpha$ under exchangeability. We use this threshold rule
unchanged; the trajectory-level content is in the score and in
the split discipline (Section 3).

\paragraph{Learned step scorers.} Process reward models score
individual steps of a reasoning trace: trained on human step
labels \citep{lightman2023verify} or on automatic process
supervision \citep{wang2024mathshepherd}, and lifted to
multi-step agent trajectories by AgentPRM
\citep{xi2025agentprm}. Trace-risk modelling without guarantees
exists as well (e.g.\ SafetyDrift's absorbing-Markov-chain
risk model over agent traces \citep{dhodapkar2026safetydrift}). None of these places a distribution-free
finite-sample guarantee over the learned score, prices
abstention against a pre-registered ceiling, or evaluates the
construction as a feasibility question with frozen data floors.
Our delta is the guarantee, the abstention pricing, and the
feasibility discipline --- not a better scorer; had the census
found eligible populations, the confirmatory signal would have
been a standard automatically-trained step scorer of exactly
this family, calibrated by CRC.

\paragraph{Corpora.} The census covers the two public
collections that survived the programme's counting screen: the
SWE-smith trajectory release (trajectories of SWE-agent driven
by Claude models on synthetic software-engineering tasks, in
three message formats) and CoderForge-Preview (OpenHands
trajectories driven by Qwen3-Coder-480B on tasks from three
sources), the latter documented by the SWE-Critix construction
\citep{swecritix2026}. Corpora excluded at the counting stage
(HAL slices, Open-SWE-Traces, earlier candidates) are excluded
for counted, documented reasons in the programme record and are
not re-litigated here.

\section{Method}

\subsection{Trajectory-level CRC (T-CRC)}

A trajectory $i$ has steps $1 \dots T_i$ and a binary outcome
(failed / passed) under its benchmark's own pass condition. A
per-step scorer $s(i, t) \in [0,1]$ estimates step-level failure
propensity; the trajectory score is the hazard aggregation
$S_{\mathrm{haz}}(i) = 1 - \prod_{t} (1 - s(i,t))$, with
$S_{\max}$ and $S_{\mathrm{mean}}$ as descriptive variants. The
confirmatory pipeline: fit the scorer on a \emph{fit} split;
pass a mechanical signal-quality gate on a disjoint \emph{gate}
split (the score must beat the trajectory-length baseline
$S_{\mathrm{len}}(i) = T_i$ decisively: AUC $\ge 0.65$,
$\Delta$AUC $\ge +0.10$, bootstrap 95\% CI for the difference
excluding 0 --- a learned signal that cannot beat ``how long
did it run'' has no business being certified); select $\lamhat$
by the CRC rule on a disjoint \emph{calibration} split; evaluate
once on a disjoint \emph{test} split, reporting risk and
abstention only as a pair. Because the gate event is measurable
with respect to the fit and gate splits, which are disjoint from
calibration and test, conditioning on it does not alter the
calibration distribution: the guarantee for the gated pipeline
is the unconditional CRC statement. A guarantee-alert rule (a
test-risk confidence interval lying entirely above $\alpha$) and
an abstention ceiling (40\% at $\alpha = 0.10$) complete the
verdict logic, frozen as an exhaustive decision table over
per-family outcomes.

\subsection{The feasibility protocol (frozen)}

Eligibility of a population is decided from a \emph{census}
alone --- artefact-level quantities computed before any signal,
gate, or split outcome exists --- under these frozen criteria:
(i) at least 800 usable trajectories (a step log and a final
evaluation both present); (ii) base failure rate in
$[0.10, 0.40]$ for the confirmatory tier, $(0.40, 0.50]$ for a
descriptive-only tier, outside $[0.10, 0.50]$ ineligible --- the
band reflects the method's operating point: below 0.10 the
class-count rule below cannot be met at feasible sizes, above
0.50 no threshold at $\alpha = 0.10$ can be non-vacuous against
a 40\% abstention ceiling; (iii) a four-way split realisation
(fit 40 / gate 20 / calibration 20 / test 20, by task group,
stratified by group outcome composition) meeting the exact
Clopper--Pearson floors $n_{\mathrm{cal}} \ge 160$ and
$n_{\mathrm{test}} \ge 160$ --- $n = 150$ gives a worst-case
95\% half-width of 0.0513, $n = 160$ gives 0.0495, and the floor
is set so the interval can resolve the $\alpha = 0.10$ decision
--- and at least 40 successes and 40 failures in each of gate,
calibration, and test.

\paragraph{Outcome binarisation as a frozen function.} The
subtlest criterion is the outcome itself. A population's
successes are defined by its benchmark's own pass condition,
applied by a rule fixed before the census sees any value: for a
boolean field, success is its true value; for a numeric field
with observed distinct-value set $V$, success is ``value $= 1$''
if every value in $V$ is exactly 0 or 1, ``value exactly $1.0$''
(a full pass; fractional values are partial passes) if every
value lies in $[0,1]$, and \emph{ineligible} if any value lies
outside $[0,1]$ or is non-numeric --- no clause may be selected
or invented after $V$ is observed. This is the criterion that
decides the CoderForge family, and we flag its strictness as a
deliberate design choice: a sentinel value that the benchmark
does not document cannot be assigned a meaning by the analyst
without converting the census into the post-hoc exercise it
exists to prevent.

\paragraph{Splits and leakage.} All splits are by task group
(every trajectory of a task in exactly one split), because the
text features of any learned scorer make task-identity leakage
a concrete failure mode when the same task appears on both
sides of a split. Task identity itself is a census item
(Section 4).

\subsection{Programme context}

The registration executed here is the programme's fourth
artefact: a v1 pilot on a HAL corpus (Section 5.4), a V2 design
retired at design review, a V3 design that found no countable
corpus at its floors, and this V4 registration over the
populations that a dedicated counting pass (artefact-level
counts and documented facts only, no outcome values) found
standing. The registration was scored before execution by a
fresh design-stage tester (design quality 7/10; venue screen
9/12, low flag, at this journal) and its execution was audited
after the fact by a second fresh tester (Section 8).

\section{Census protocol}

Each population was censused by streaming its artefact files at
the revision counted in the prior screen (CoderForge-Preview at
revision \texttt{060fca96}; SWE-smith trajectories at
\texttt{08e109b4}), one file at a time with per-file checkpoints,
retaining only aggregates. Per population the census records:
rows in file vs.\ usable rows; parse failures (zero in all nine
populations); the outcome field's distinct values, verbatim,
with counts; the task-identifier field as found and the number
of distinct tasks; median and mean trajectory length in steps;
and step coverage. SWE-smith populations are (message format
$\times$ model) cells of one release, filtered on its model
field; its format cells share one task pool --- the xml and
ticks cells for the same model are the same trajectories in
different formats, with identical outcome counts --- and its
cells are therefore never treated as independent corpora.

\paragraph{Task identity.} SWE-smith carries \texttt{instance\_id}
as a column. CoderForge carries no task-id column; schema
identification at census found the task identity in the
artefact's \texttt{image} field (the per-task environment image;
mechanically verified on one inspected file of 148 rows:
\texttt{trajectory\_id} minus its run suffix equals the image
string on 148/148 rows). The final distinct counts --- CF-1
4{,}492 images over 77{,}169 trajectories, CF-2 444 over
148{,}001, CF-3 4{,}274 over 32{,}964 --- show these are
per-repository environment images of varying granularity, not
per-instance identifiers (near-instance only for CF-3).
Grouping by image would have been task-conservative (no task
can straddle splits when grouped by a coarser key) but coarser
than task grouping. An independent adjudication of this
identification, its timing (recorded before any eligibility
decision), and its consequences is part of the programme record
and is summarised in Section 8: it changes no outcome, since
CoderForge's eligibility fails on the outcome encoding under
any reading of the task-identity question.

\section{Results}

\subsection{The census}

Table~\ref{tab:census} is the result. Of 332{,}499 censused
trajectories (the sum over the nine registered populations;
the two families partially share the SWE-smith task universe,
and SS-2/SS-3 duplicate one another's trajectories across
formats), \textbf{no population is eligible} under the frozen
criteria, and no population reached the split realisation: the
sample draw, the split-formation checks, and eligibility leg
(iii) were never evaluated for any population, because every
population failed at the outcome stage --- leg (ii) for
SWE-smith, the binarisation rule for CoderForge. The programme
verdict under the registration's own decision table (both
families without an eligible population) is the pre-registered
\textbf{H-DATA} outcome: the confirmatory programme does not
run, and the census is the finding.

\begin{table}[t]
\centering\small
\begin{tabular}{lrrrl}
\toprule
Population & Usable & Tasks & Binding quantity & Outcome\\
\midrule
CF-1 (SWE-Rebench source) & 77{,}169 & 4{,}492 & $-$1.0 on 255 rows (0.33\%) & ineligible\\
CF-2 (SWE-Smith source) & 148{,}001 & 444 & $-$1.0 on 367 rows (0.25\%) & ineligible\\
CF-3 (R2E-Gym source) & 32{,}964 & 4{,}274 & $-$1.0 on 278 rows (0.84\%) & ineligible\\
SS-1 (tool, Claude 3.7) & 17{,}715 & 9{,}143 & base failure 0.6102 & ineligible\\
SS-2 (xml, Claude 3.7) & 18{,}457 & 11{,}824 & base failure 0.5447 & ineligible\\
SS-3 (ticks, Claude 3.7) & 18{,}457 & 11{,}824 & base failure 0.5447 & ineligible\\
SS-4 (tool, Claude 3.5) & 5{,}751 & 1{,}965 & base failure 0.5921 & ineligible\\
SS-5 (xml, Claude 3.5) & 7{,}091 & 2{,}223 & base failure 0.5875 & ineligible\\
SS-6 (ticks, Claude 3.5) & 6{,}894 & 2{,}222 & base failure 0.5863 & ineligible\\
\bottomrule
\end{tabular}
\caption{The feasibility census. ``Binding quantity'' is the
first frozen criterion each population fails: for CoderForge,
the outcome value set contains $-1.0$, which the frozen
binarisation function's third clause makes census-ineligible;
for SWE-smith, the base failure rate exceeds the 0.50 ceiling
of eligibility leg (ii) (leg (i), usable $\ge 800$, passes for
all six). Usable counts reproduce the prior counting screen
exactly; parse failures are zero in all nine populations.
Median trajectory lengths: CF 48--59 steps, SS 15--30 steps.}
\label{tab:census}
\end{table}

\subsection{Barrier 1: base-rate concentration (SWE-smith)}

Every Claude cell of the SWE-smith release fails at a base
failure rate between 0.545 and 0.610. This is not a threshold
artefact of the census: the populations are large (5{,}751 to
18{,}457 usable trajectories), the outcome field is an
unambiguous boolean (\texttt{resolved}), and the rates are what
they are --- SWE-agent driven by these models resolves roughly
two in five of these synthetic tasks. A confirmatory method
whose registered operating point is $\alpha = 0.10$ with a 40\%
abstention ceiling cannot certify such populations without
abstaining on nearly everything, which the anti-vacuity
discipline of the protocol exists to refuse. The band's ceiling
is doing intended work here.

\subsection{Barrier 2: an uninterpretable outcome value (CoderForge)}

CoderForge's \texttt{reward} is a float taking values 1.0
(89{,}501 $+$ 44{,}739 $+$ 20{,}904 rows across the three
sources), 0.0, and $-1.0$ --- 255, 367, and 278 rows in CF-1,
CF-2, CF-3 respectively. No definition of the $-1.0$ value
appears in the artefact documentation or in the counting and
census records we inspected; plausible readings exist (an
execution or evaluation failure sentinel is the obvious one),
but the frozen binarisation has no clause for a value outside
$[0,1]$, by design, and inventing one at census time is exactly
the move the pre-registration prohibits. All three populations
are therefore ineligible on the outcome encoding.

\paragraph{Exploratory counterfactuals (unregistered).} To make
the cost of that discipline visible --- and only for that
purpose --- we record what two unregistered readings of $-1.0$
would have given. Excluding the $-1.0$ rows, the CF failure
shares are CF-1 0.4183, CF-2 0.3938, CF-3 0.3605. Counting
$-1.0$ as failure gives CF-1 0.4202, CF-2 0.3953, CF-3 0.3659.
The band assignment is the same under both readings: CF-2 and
CF-3 would have sat inside the confirmatory band, CF-1 in the
descriptive band (it did not attain descriptive-tier status,
which is a registered status, and we do not describe it as a
descriptive-tier population). For CF-2 and CF-3 the binding
barrier is thus the outcome encoding, not the band; that is a
statement about documentation, obtainable only because the
registered rule refused to make it silently.

\subsection{Machinery evidence}

The bottleneck is not the guarantee machinery. The programme's
validation harness (48 checks, all passing) includes the
classical anchors --- single-turn CRC reproducing its nominal
risk (mean test risk 0.1000 at $\alpha = 0.10$), planted
-dependence synthetics controlled by the hazard score (mean
risk 0.0154) --- and a synthetic prefix anchor for the full
learned pipeline eventually intended for the confirmatory runs:
on synthetic traces with a planted per-step error signal, the
frozen feature map, stacked heads, and hazard aggregation
recover the signal (gate AUC 0.995) and CRC controls the risk
at $\alpha = 0.10$ (test risk 0.080), at an abstention of
0.420 --- above the 40\% ceiling, so the anchor demonstrates
risk control, not a confirmatory pass, and we report it as
such. What the census establishes is that this machinery
currently has no eligible natural population to run on.

\subsection{The prior pilot, in its place}

Under the programme's separate v1 registration, one HAL corpus
(tau-bench airline) was run end-to-end with a frozen heuristic
signal ladder: base failure 0.84, calibration AUC 0.482,
threshold at abstain-all for every $\alpha$, abstention 100\%
at $\alpha = 0.10$ --- a falsifier-direction result for that
signal class on that corpus. It is v1 evidence, cited for the
programme's history; the present verdict does not rest on it,
and its corpus's 0.84 base rate would itself be ineligible
under the present band, which post-dates it and has no
retroactive effect on its recorded verdict.

\section{Discussion}

\paragraph{What the result says.} Under criteria frozen before
outcomes were seen, the two public software-engineering agent
trajectory collections that survive a counted screen cannot
support confirmatory trajectory-level conformal risk control:
one family's populations all sit above the base-rate ceiling,
and the other's outcome field cannot be binarised without
inventing a meaning for an undocumented value. Both barriers
are properties of the data landscape, not of the guarantee:
the same protocol, unchanged, would have run on any population
that met it, and its machinery anchors pass. Pre-registration
is what makes this a result rather than an excuse: the H-DATA
outcome was one of the registration's pre-declared verdicts,
reached by its own decision table, with floors and bands that
were never adjusted after the data arrived.

\paragraph{What it does not say.} It does not say trajectory
-level CRC fails on agent traces: no eligible population was
ever run, and the one historical run (Section 5.4) tested a
heuristic ladder, not the learned signal. It does not say the
band is ``right'' in any absolute sense: the band encodes an
operating point ($\alpha = 0.10$, 40\% abstention ceiling), and
a programme targeting a different operating point would draw
it elsewhere --- but would have to draw it in advance, as this
one did. And it does not generalise beyond its census: two
artefact families, one task domain, nine populations.

\paragraph{The cost of the discipline, made visible.} The
exploratory counterfactuals of Section 5.2 are the honest
price tag of the frozen binarisation: two populations that a
post-hoc analyst would very likely have run (failure shares
0.36--0.39 under either plausible reading of $-1.0$) were
turned away over 0.25--0.84\% of rows whose meaning the
artefact does not document. We judge that price correctly
paid --- a guarantee certified on an invented outcome
definition is not a guarantee --- but it is a price, and
corpus publishers can reduce it to zero for everyone
(Section 7).

\section{Recommendations}

We label these as recommendations; each is tied to the census
fact that motivates it, and none is an experimental finding.

\begin{enumerate}
\item \textbf{Document sentinel outcome values, or encode
execution failure separately from task failure.} An undefined
$-1.0$ reward value, present in all three CoderForge sources
at 0.25--0.84\% of rows, is the sole reason three countable
populations could not be binarised under a pre-registered
rule.
\item \textbf{Ship the benchmark's pass/fail outcome as a
boolean field.} SWE-smith's \texttt{resolved} boolean was the
one outcome encoding in this census that required no
interpretation whatsoever.
\item \textbf{Carry task/instance identifiers as a column.}
SWE-smith's \texttt{instance\_id} column made task-grouped
splitting unambiguous; CoderForge has no task-id column, and
its environment-image field is a per-repository identifier
whose granularity varies by source (444 to 4{,}492 distinct
values over 33{,}000--148{,}000 trajectories).
\end{enumerate}

\section{Limitations}

The census covers two artefact families in one task domain
(software-engineering agents); other domains' corpora were
excluded at the programme's counting stage for documented
reasons and are outside this paper's claims. The CoderForge
task-identity identification (Section 4) rests on mechanical
verification of one inspected file and an adjudicated
interpretation; it is outcome-immaterial here but would matter
in any future run on this artefact, where per-instance
identifiers would have to come from the publisher or from the
trajectories' own content under a new registration. The
exploratory rates of Section 5.2 are unregistered
counterfactuals, reported for interpretation; neither reading
of $-1.0$ is endorsed. The synthetic machinery evidence
(Section 5.3) uses planted signals and demonstrates control,
not confirmatory performance. Finally, the eligibility band
and floors are design choices: we argue for them from the
method's operating point, and other pre-registered choices are
possible --- what is not optional, on this paper's argument,
is making them in advance.

\paragraph{Falsifier bookkeeping.} The programme sits inside a
larger area whose pre-registered falsifier requires
non-vacuous trajectory-level control (at most 40\% abstention
at target risk) on at least two of three trace corpora: a HAL
task family, own-lab traces, and METR runs. The HAL family
carries the v1 falsifier-direction result; the other two are
untested; the falsifier has therefore neither fired nor been
passed. The present census's families are none of those three,
and its results count neither for nor against that falsifier.
H-DATA is this programme's verdict under its own registration;
any area-level kill decision is a separate, recorded act by
the programme's owner, not a consequence this paper asserts.

\paragraph{Reproducibility.} Every number in Table~\ref{tab:census}
recomputes from the census records produced by the programme's
census code over the two artefacts at the revisions named in
Section 4. On freeze integrity, precisely: the programme's
earlier registration chains (V2, V3) and the present
registration chain's non-registration files verify against
their recorded hashes today; the present registration's Part I
was hash-verified at its pre-execution design-stage check, and
its Part II execution record was appended under the
registration's own rule. The code, registrations, census
records, and tester reports are deposited in a public
repository: \url{https://github.com/ranadhir19/trajectory-conformal-census}.

\begin{thebibliography}{9}
\small
\bibitem[Angelopoulos et~al.(2022)]{angelopoulos2022conformal}
Angelopoulos, A.~N., Bates, S., Fisch, A., Lei, L., and
Schuster, T. (2022). Conformal risk control.
\emph{arXiv:2208.02814}.

\bibitem[Lightman et~al.(2023)]{lightman2023verify}
Lightman, H., Kosaraju, V., Burda, Y., Edwards, H., Baker, B.,
Lee, T., Leike, J., Schulman, J., Sutskever, I., and Cobbe, K.
(2023). Let's verify step by step. \emph{arXiv:2305.20050}.

\bibitem[Wang et~al.(2024)]{wang2024mathshepherd}
Wang, P., Li, L., Shao, Z., Xu, R.~X., Dai, D., Li, Y., Chen,
D., Wu, Y., and Sui, Z. (2024). Math-Shepherd: Verify and
reinforce LLMs step-by-step without human annotations.
\emph{arXiv:2312.08935}.

\bibitem[Xi et~al.(2025)]{xi2025agentprm}
Xi, Z., Liao, C., Li, G., Yang, Y., Chen, W., Zhang, Z., Wang,
B., Jin, S., Zhou, Y., Guan, J., Wu, W., Ji, T., Gui, T.,
Zhang, Q., and Huang, X. (2025). AgentPRM: Process reward
models for LLM agents via step-wise promise and progress.
\emph{arXiv:2511.08325}.

\bibitem[Dhodapkar and Pishori(2026)]{dhodapkar2026safetydrift}
Dhodapkar, A. and Pishori, F. (2026). SafetyDrift: Predicting
when AI agents cross the line before they actually do.
\emph{arXiv:2603.27148}.

\bibitem[Hu and Chong(2026)]{swecritix2026}
Hu, J. and Chong, C.~Y. (2026). SWE-Critix: Post-training
generalist verifiers for assessing coding-agent trajectories.
\emph{arXiv:2610.13356}.
\end{thebibliography}

\end{document}
